% Vyžaduje balíčky: graphicx, float. Obrázky: analyza/analyticka_cast/obrazky/
% (\graphicspath{{analyza/analyticka_cast/obrazky/}}).
\section{Prostorový rozměr}

Prostorové rozložení objektů ukazují mapy vytvořené z definičních bodů
objektů, které má vyplněné 99,9\,\% aktuálních objektů; obrys budovy má jen
97,8\,\%. Hustota a průměrné hodnoty jsou proto počítány v pravidelné
šestiúhelníkové síti o velikosti buňky přibližně 250\,m, ne na obrysech.
Mapy jsou bez podkladové mapy a hranic obcí, severní směr je nahoře.

\subsection{Hustota zástavby}

Zástavba tvoří v Brně souvislé jádro, od kterého vybíhají úzké pásy podél
údolí a komunikací (obrázek~\ref{fig:hustota}). Okolní obce, například Kuřim,
Šlapanice, Modřice, Bílovice nad Svitavou nebo Střelice, vytvářejí
oddělené shluky.

\begin{figure}[H]
\centering
\includegraphics[width=\textwidth]{23_mapa_hustota}
\caption{Hustota stavebních objektů v šestiúhelníkové síti}
\label{fig:hustota}
\end{figure}

\subsection{Podlažnost a stáří zástavby}

Podlažnost a stáří zástavby se v prostoru výrazně liší. Průměrně mají objekty
v Brně 2,5 podlaží, v okolních obcích 1,75. Nejvyšší zástavba je v centru
a v oblastech velkých sídlišť, kde průměr překračuje čtyři podlaží
(obrázek~\ref{fig:podlaznost-mapa}). Okolí je převážně nízké, s průměrem
jedno až dvě podlaží.

\begin{figure}[H]
\centering
\includegraphics[width=\textwidth]{24_mapa_podlaznost}
\caption{Průměrná podlažnost objektů v buňkách sítě (buňky s méně než 10 objekty se nevykreslují)}
\label{fig:podlaznost-mapa}
\end{figure}

Stejně se liší stáří (obrázek~\ref{fig:rok-mapa}). Medián roku dokončení je
v Brně 1980, v okolních obcích 2001. Nejstarší zástavba leží v jádru města,
nejnovější je v Kuřimi a Modřicích (medián 2004). U starších objektů je
nutné pamatovat na to, že rok dokončení je často jen hranicí období
(viz kapitola o stáří zástavby).

\begin{figure}[H]
\centering
\includegraphics[width=\textwidth]{25_mapa_rok_dokonceni}
\caption{Medián roku dokončení v buňkách sítě (buňky s méně než 10 objekty se nevykreslují)}
\label{fig:rok-mapa}
\end{figure}

\subsection{Nová výstavba}

V letech 2015--2024 bylo dokončeno 8\,197 objektů, z toho 4\,689 v Brně
(obrázek~\ref{fig:nova-mapa}). Výstavba je rozptýlená po celém území, výrazné
shluky vznikají na okrajích Brna a v okolních obcích. Vzhledem k množství
stávající zástavby je nová výstavba v okolních obcích výrazně intenzivnější:
na 1\,000 stávajících objektů připadá 125 nově dokončených v okolních
obcích a 63 v Brně (obrázek~\ref{fig:nova-obce}). Nejvíce stavělo
Moravany (163), Střelice (159) a Popůvky (153), nejméně Šlapanice (62).

\begin{figure}[H]
\centering
\includegraphics[width=\textwidth]{26_mapa_nova_vystavba}
\caption{Objekty dokončené v letech 2015--2024 (oranžově) nad pozadím všech objektů}
\label{fig:nova-mapa}
\end{figure}

\begin{figure}[H]
\centering
\includegraphics[width=0.9\textwidth]{27_nova_vystavba_podle_obci}
\caption{Nová výstavba v letech 2015--2024 v největších obcích na 1\,000 stávajících objektů}
\label{fig:nova-obce}
\end{figure}
